\begin{table}[htbp]
  \centering
  \small
  \caption{Matriz de Probabilidades de Transición, Distribución Estacionaria y Tiempos de Residencia}
  \label{tab:transicion_markov}
  \label{tab:markov}
  \begin{tabular}{lcccccc}
    \toprule
    Estado & \multicolumn{3}{c}{Probabilidad de Transición hacia Estado Destino} & Distr. Estacionaria & T. Residencia \\
    \cmidrule(lr){2-4}
    Origen ($S_t$) & Positivo ($S_{t+1}$) & Negativo ($S_{t+1}$) & Neutro ($S_{t+1}$) & ($\pi_j$) & (turnos) \\
    \midrule
    Positivo & 0,160 & 0,064 & 0,776 & 0,115 (11,5\%) & 1,19 \\
    Negativo & 0,102 & 0,075 & 0,823 & 0,054 (5,4\%)  & 1,08 \\
    Neutro   & 0,109 & 0,051 & 0,839 & 0,831 (83,1\%) & 6,22 \\
    \bottomrule
  \end{tabular}
  \vspace{4pt}
  \begin{minipage}{0.94\textwidth}
    \footnotesize
    \textit{Nota.} Tiempo de recuperación ante episodios negativos: $\tau_{rec} = 1,08$ turnos. El análisis estocástico demuestra que el estado neutro/factual actúa como un atractor cuasi-absorbente dentro de la dinámica de conversación universitaria, absorbiendo más del 77\% de las transiciones salientes de los estados polares.
  \end{minipage}
\end{table}
